\documentclass[12pt]{article}
\usepackage[utf8]{inputenc}
\usepackage[T1]{fontenc}
\usepackage{hyperref}
\title{Vod Language Resources for K-12 Curriculum Development}
\author{Compiled by Claude Sonnet 5 (claude-sonnet-5) for Andr\'as Kornai}
\date{}
\begin{document}
\maketitle

\section{Introduction}
Vod (Votic, Vadja; ISO 639-3 vot) is a critically small Uralic language of the Ingria region, with only elderly speakers. Most sources are in Estonian, Russian or Finnish.

\section{Primary Educational Resources}
An IDIL 2022 teaching kit for Vod and Izhora (27 lessons, short dictionary of over 200 words, in Russian) was mentioned; the hosting platform has been inactive since February 2025: \url{https://idil2022-2032.org/?p=51350}

\section{Linguistic and Reference Materials}
\begin{itemize}
\item Electronic Morphological Dictionary of Votic (Rozhanskiy and Markus; about 1,800 entries, Russian and English, paradigms, about 5,800 audio examples): \url{https://ingrian.org/votic-dictionary}
\item Votian orthography (Wikiversity): \url{https://en.wikiversity.org/wiki/Votian_orthography}
\item Wiktionary Votic entry guidelines: \url{https://en.wiktionary.org/wiki/Wiktionary:AVOT}
\item Printed works known from citations: Adler and Leppik, \emph{Vadja keele s\~onaraamat}; Ariste, \emph{A Grammar of the Votic Language} (1968). No online links to these were verified.
\end{itemize}

\section{Religious and Community Resources}
None found.

\section{Digital and Audio Resources}
The electronic dictionary above includes recordings from last first-language speakers. Endangered Languages Project listing (three media items): \url{https://endangeredlanguages.com/node/38612/resources}

\section{Curriculum Development Context}
Reference and audio documentation exists, and a short teaching kit exists, but no school curriculum was found. Revitalization would build on the dictionary audio and the IDIL kit.
\end{document}